\chapter{Introduction}

Medicinal and aromatic plants (MAPs) represent an important group of crops widely used in traditional medicine, food production, cosmetics, and pharmaceutical industries due to their rich content of bioactive compounds \textcite{Ahmad2020}. These plants produce secondary metabolites such as essential oils, phenolic compounds, and terpenoids that exhibit antimicrobial, antioxidant, and anti-inflammatory properties \textcite{Ahmad2020}. Because of the growing global demand for natural products and plant-derived medicines, the cultivation of medicinal and aromatic plants has increased significantly in recent decades \textcite{Ahmad2020}.
\\\\
Among medicinal and aromatic plants, species belonging to the genus \textit{Mentha} (mint) are widely cultivated and economically important herbs from the Lamiaceae family \textcite{Kalemba2019}. Mint species are valued for their high biomass production and essential oil content, which makes them important raw materials for the food, cosmetic, and pharmaceutical industries \textcite{Zheljazkov2010}. The essential oils extracted from mint plants contain several biologically active compounds, including menthol, menthone, and carvone, which are responsible for their characteristic aroma, flavor, and therapeutic properties \textcite{Ahmad2020}. These compounds are widely used as flavoring agents in beverages, confectionery products, oral hygiene products, and medicinal preparations \textcite{Ahmad2020}.
\\\\
Among the numerous species within the genus, \textit{Mentha piperita} (peppermint), \textit{Mentha spicata} (spearmint), and \textit{Mentha crispata} are considered particularly important due to their high essential oil yield and extensive commercial cultivation \textcite{Kalemba2019}. Peppermint is a hybrid species derived from \textit{Mentha aquatica} and \textit{Mentha spicata} and is widely cultivated for its oil rich in menthol and menthone, compounds that contribute to its cooling sensation and medicinal effects \textcite{Ahmad2020}. Spearmint is characterized by the presence of carvone as a major component of its essential oil and is widely used in culinary products and herbal formulations \textcite{Zheljazkov2010}. Due to their chemical composition and biological activities, these mint species have gained considerable importance in agricultural production and the global herbal market \textcite{Kalemba2019}.
\\\\
In modern horticultural systems, medicinal and aromatic plants are often cultivated under greenhouse or controlled conditions to ensure consistent yield and quality \textcite{Mashamaite2024}. In such systems, the selection of appropriate growing media is essential for plant growth and development. Peat moss has long been one of the most commonly used substrate components because of its favorable properties, including high water-holding capacity, low bulk density, and good aeration \textcite{Mashamaite2024}. However, peat extraction has raised serious environmental concerns because peatlands function as important carbon reservoirs and their degradation can release large amounts of stored carbon into the atmosphere \textcite{Mashamaite2024}. Consequently, the horticultural sector has increasingly focused on identifying sustainable alternatives to peat-based growing media.
\\\\
Biochar has recently emerged as a promising alternative substrate component in sustainable agriculture and horticulture \textcite{Mashamaite2024}. Biochar is a carbon-rich material produced through the pyrolysis of organic biomass under limited oxygen conditions \textcite{Mashamaite2024}. Due to its porous structure, large surface area, and high cation exchange capacity, biochar can improve the physical and chemical properties of growing media by enhancing water retention, nutrient availability, and microbial activity \textcite{Mashamaite2024}. In addition to improving soil properties, biochar has also been reported to alleviate environmental stresses in plants, including salinity stress \textcite{Farhangi2017}.
\\\\
Salinity stress is considered one of the most important abiotic stresses limiting plant growth and agricultural productivity worldwide \textcite{Volkov2014}. High concentrations of soluble salts in the root zone can reduce water uptake by plants, cause ionic toxicity, and disrupt nutrient balance \textcite{Volkov2014}. The accumulation of sodium ions in plant tissues can interfere with metabolic processes, enzyme activity, and photosynthesis, ultimately leading to reduced plant growth and biomass production \textcite{Volkov2014}. Medicinal and aromatic plants, including mint species, are particularly sensitive to salinity because environmental stress can affect both plant growth and the synthesis of secondary metabolites responsible for their medicinal value \textcite{Kalemba2019}.
\\\\
Recent studies have suggested that biochar application can help mitigate the negative effects of salinity by improving soil properties and reducing sodium availability in the root zone \textcite{Farhangi2017}. Biochar amendments can enhance plant tolerance to salt stress by improving water retention, increasing nutrient availability, and promoting beneficial microbial activity in the soil \textcite{Mashamaite2024}. However, the response of plants to biochar amendments varies depending on plant species, biochar characteristics, and environmental conditions \textcite{Mashamaite2024}. In addition, limited research has examined how different mint species respond to salinity stress when grown in biochar-amended substrates.
\\\\
Therefore, understanding the interaction between growing media composition and salinity stress is important for developing sustainable cultivation strategies for medicinal and aromatic plants. In this context, the present study evaluates the effects of biochar as a potential alternative to peat-based growing media on the salt stress responses of three mint species: \textit{Mentha piperita}, \textit{Mentha spicata}, and \textit{Mentha crispata}. The study aims to assess plant growth and physiological responses under saline conditions and to determine whether biochar can improve plant tolerance to salt stress while serving as an environmentally friendly substrate component.